\section{Introduction}

Advanced Persistent Threat (APT) campaigns unfold as multi-stage operations that may span
days or months before an objective is reached. Recent industry reporting places the global
median dwell time at 14 days in 2025, with cyber-espionage intrusions found considerably
later\cite{mandiant2026mtrends}, attributes a large share of intrusions to credential abuse
and vulnerability exploitation\cite{verizon2025dbir}, and records substantial business
disruption once a breach occurs\cite{ibm2024costdatabreach}. Such figures do not convey the
shape of the defense. Individual actions or stages may have been detected or disrupted even
when the campaign ultimately succeeded, and a campaign that failed may have been stopped by a
single control acting late. Recording only whether a campaign succeeded therefore hides where
defensive capability held and where it was insufficient.

Organizations respond by deploying layered stacks, but deployment is not capability. The
components of such a stack emit evidence with different meanings: one removes a precondition
so that an action cannot be attempted, another interrupts the action while it runs, a third
records it without preventing it, and a fourth does not observe it. Counting products, or
counting alerts, flattens these distinctions. An evaluation useful to a defender must instead
be comparable across heterogeneous stacks, aware of where in a campaign an outcome occurred,
and traceable from an aggregate figure back to a control that can be repaired.

Existing approaches supply parts of this picture. ATT\&CK provides a shared vocabulary for
adversary behavior\cite{mitreAttack} and public product evaluations report responses step by
step\cite{mitreAttackEvaluations}, yet technique coverage has been shown empirically to be a
weak proxy for defensive capability\cite{virkud2024endpoint}. Provenance-based detection and
attack reconstruction recover the causal origin of observed malicious
activity\cite{jiang2025orthrus,goyal2024rcaid}, and control catalogs relate adversary
techniques to countermeasures\cite{mitreCTIDMappings,kaloroumakis2021d3fend}. These outputs
characterize attacker behavior, causal evidence, or candidate countermeasures, but they are
not designed to provide the stable victim-side control attribution required by our
repair-oriented scoring task.

A second line represents how an intrusion advances. Attack graphs and reachability models
reason about attack paths and their
composition\cite{sheyner2002attackgraphs,wang2007measuring}, while severity and system-level
risk models quantify exposure along a different axis\cite{mell2007common,erola2022system};
step-level results have also been re-linked into causal chains to recover campaign
structure\cite{shen2024decoding}. Campaign progress and the control-point domain associated
with an action nevertheless remain orthogonal dimensions: the same control domain can be
implicated at multiple phases, while actions within one phase can map to different control
domains. Collapsing them onto a single axis makes a score harder to act on, not easier.

A third line aggregates observations into an overall figure. Security-metric and
evaluation-validity work reports coverage, detection rates, and paired effectiveness
quantities\cite{iannacone2020quantifiable,zaffarano2015mtd,liu2023defia}, and cautions that
aggregate security numbers are sensitive to how they are
constructed\cite{verendel2009quantified,vanderkouwe2019benchmarkingflaws}. A uniform average
treats action outcomes as interchangeable. Actions within one stage are not: sometimes any
one of them suffices for the attacker, sometimes all are required, and sometimes they form a
loose set. An average cannot represent these alternative, conjunctive, and clustered stage
semantics, and it obscures which actions and stages account for the resulting score loss.

The three limitations above point to a common missing object: an evaluation unit that keeps
attacker behavior, a benchmark-side repair reference, and the evaluated system's response
aligned without conflating them. We therefore pair an in-scope benchmark action with a
reference primary control-point attribution fixed during benchmark construction, and with the
outcome observed when a defense system under test is exercised against that action. The
attribution names one victim-side control point -- a mechanism, configuration, policy, or
operational process -- and is a benchmark-side reference, assigned under a documented
procedure before any system is evaluated. The observed outcome, by contrast, records what the
evaluated system did when the action was exercised. This separation matters because an
attacker-side technique label does not by itself determine a victim-side repair target:
actions sharing the same technique label may require different control-point attributions
depending on their context. Attacker-side and victim-side attribution answer different
questions, and a framework aimed at repair must preserve the second without discarding the
first.

Building an evaluation on this unit raises three methodological challenges. First, the
attribution must remain stable across heterogeneous actions while staying repair-relevant: a
schema too coarse cannot separate distinct repairs, and one too fine cannot be applied
consistently. Second, campaign phase and control-point domain must be carried as separate
dimensions, together with the confidence of each assignment, so that a
downstream figure can be interpreted rather than merely read. Third, action-level outcomes
must be aggregated under the semantics that actually hold inside a stage, and the aggregate
must remain decomposable afterwards.

We address these with a three-step framework. The first step constructs the reference
attribution for benchmark actions, recording a primary control-point reference together with
its confidence under the benchmark schema. Secondary annotations, when present, remain
contextual metadata rather than additional primary scoring targets. The second step enriches
the attributed action into a scoring-ready record, adding the campaign phase at which the
action occurs, the defense product categories expected to address the attributed control
point, and the aggregation semantics governing its stage; phase is carried independently of
the attribution, so that where an action occurs in the campaign remains separate from which
control-point domain it is associated with. The third step aligns observations from a system
under test to benchmark actions, normalizes them onto a common outcome vocabulary of
\texttt{PREVENTED}, \texttt{BLOCKED}, \texttt{DETECTED}, and \texttt{MISSED}, and aggregates
actions to stages and stages to playbooks while retaining breakdowns by phase, control-point
domain, and product category. The current benchmark instance draws on 53 APT playbooks and
1,796 authoritative raw action occurrences; the benchmark presently defines 1,773 active
scoring opportunities, and the control-point schema presently comprises nine top-level domains
and 97 active leaves.

The main contributions of this paper are as follows:
\begin{itemize}
    \item \textbf{An evaluation representation for repair-oriented APT defense scoring.} We
    define the evaluation record so that it extends an attacker-side action description by
    binding to it a reference primary control-point attribution and an observed defense
    outcome, keeping the benchmark-side attribution separate from the behavior of the system
    under test. The attack action remains part of the representation; what the representation
    adds is a repair-oriented victim-side dimension alongside the observed result.

    \item \textbf{A scoring-ready, stage-aware benchmark representation.} We describe a
    construction process that enriches each attributed action with its campaign phase, the
    expected defense product categories, confidence, and explicit stage
    aggregation semantics, keeping campaign progress and victim-side control attribution as
    separate dimensions of one record.

    \item \textbf{A stage-aware aggregation and reporting model.} We separate an action's
    outcome from the weight of the stage containing it, aggregate outcomes under each stage's
    recorded OR, AND, or Cluster semantics rather than a uniform average, compose action to
    stage to playbook scores, and preserve decomposition by action, control point, and phase.
\end{itemize}
